% GENERATED FILE. Edit canonical lessons, then run npm run generate:books.
% canonical-source-hash: fnv1a64:17e970418d0f61dc
% canonical-lessons: FR-C314-deceder, FR-C314-enterrer, FR-C314-enterrement, FR-C314-funerailles, FR-C314-cimetiere

\chapter{Les funérailles}
\label{ch:fr-les-funerailles}
\hlchaptermodality{314}

\begin{chapteropening}
\textbf{By the end of this chapter:} \emph{I can talk about a death and a funeral: passing away, burying, the burial, the funeral and the cemetery.}

Complete the last lesson of chapter 314: I can talk about a death and a funeral: passing away, burying, the burial, the funeral and the cemetery.
\end{chapteropening}

\section[décéder]{décéder — to pass away, to die (formal)}

\label{lesson:FR-C314-deceder}



\begin{quote}
Before the new one: say the French for to take into hospital, then the French for the death.
\end{quote}

\subsection*{You'll want to know: décéder}
\textbf{décéder} — \textquotedblleft{}to pass away, to die (formal)\textquotedblright{}.

\begin{quote}
Il est décédé à l'âge de quatre-vingts ans.
\end{quote}

\begin{quote}
\emph{He passed away at the age of eighty.}
\end{quote}

Takes \emph{être}. Official language; \emph{mourir} in everyday speech.

\begin{cousinweb}
\emph{décéder} is a learned borrowing from Latin \emph{dēcēdere} \textquotedblleft{}to go away, to die\textquotedblright{} (EtymDB), from \emph{dē} and \emph{cēdere} \textquotedblleft{}to go\textquotedblright{}.
\end{cousinweb}

\begin{grammarlens}[title={What you've built}]
One more word for telling, explaining and describing at B1.
\end{grammarlens}

\subsection*{Your turn}
\begin{itemize}
\raggedright
  \item \emph{Say it:} \emph{décéder}
  \item \emph{Say it:} \emph{décéder}, once more
  \item \emph{Say it:} a sentence of your own with \emph{décéder}
  \item \emph{Recall:} say the French for to take into hospital, then the French for the death, then say \emph{décéder} again
\end{itemize}

\begin{tcolorbox}[breakable,colback=teal!4,colframe=teal!35!black,title={Before you move on}]
What does décéder mean? (\textquotedblleft{}To pass away, to die (formal)\textquotedblright{}.) Say it once more.
\end{tcolorbox}
\section[enterrer]{enterrer — to bury}

\label{lesson:FR-C314-enterrer}



\begin{quote}
Before the new one: say the French for the death, then the French for to pass away.
\end{quote}

\subsection*{You'll want to know: enterrer}
\textbf{enterrer} — \textquotedblleft{}to bury\textquotedblright{}.

\begin{quote}
Ils ont enterré leur grand-père près de l'église.
\end{quote}

\begin{quote}
\emph{They buried their grandfather near the church.}
\end{quote}

Figurative: \emph{enterrer un projet} — to drop a plan.

\begin{cousinweb}
\emph{enterrer} is inherited from Vulgar Latin \emph{interrāre} \textquotedblleft{}to put in the earth\textquotedblright{} (EtymDB), built on Latin \emph{terra} \textquotedblleft{}earth\textquotedblright{}.
\end{cousinweb}

\begin{grammarlens}[title={What you've built}]
One more word for telling, explaining and describing at B1.
\end{grammarlens}

\subsection*{Your turn}
\begin{itemize}
\raggedright
  \item \emph{Say it:} \emph{enterrer}
  \item \emph{Say it:} \emph{enterrer}, once more
  \item \emph{Say it:} a sentence of your own with \emph{enterrer}
  \item \emph{Recall:} say the French for the death, then the French for to pass away, then say \emph{enterrer} again
\end{itemize}

\begin{tcolorbox}[breakable,colback=teal!4,colframe=teal!35!black,title={Before you move on}]
What does enterrer mean? (\textquotedblleft{}To bury\textquotedblright{}.) Say it once more.
\end{tcolorbox}
\section[l'enterrement]{l'enterrement — the burial, the funeral}

\label{lesson:FR-C314-enterrement}



\begin{quote}
Before the new one: say the French for to pass away, then the French for to bury.
\end{quote}

\subsection*{You'll want to know: l'enterrement}
\textbf{l'enterrement} — \textquotedblleft{}the burial, the funeral\textquotedblright{}.

\begin{quote}
Beaucoup de gens sont venus à l'enterrement.
\end{quote}

\begin{quote}
\emph{Many people came to the funeral.}
\end{quote}

The everyday word for a funeral.

\begin{cousinweb}
\emph{enterrement} is built on \emph{enterrer} \textquotedblleft{}to bury\textquotedblright{}, which is inherited from Vulgar Latin \emph{interrāre}, built on \emph{terra} \textquotedblleft{}earth\textquotedblright{} (EtymDB).
\end{cousinweb}

\begin{grammarlens}[title={What you've built}]
One more word for telling, explaining and describing at B1.
\end{grammarlens}

\subsection*{Your turn}
\begin{itemize}
\raggedright
  \item \emph{Say it:} \emph{l'enterrement}
  \item \emph{Say it:} \emph{l'enterrement}, once more
  \item \emph{Say it:} a sentence of your own with \emph{l'enterrement}
  \item \emph{Recall:} say the French for to pass away, then the French for to bury, then say \emph{l'enterrement} again
\end{itemize}

\begin{tcolorbox}[breakable,colback=teal!4,colframe=teal!35!black,title={Before you move on}]
What does l'enterrement mean? (\textquotedblleft{}The burial, the funeral\textquotedblright{}.) Say it once more.
\end{tcolorbox}
\section[les funérailles]{les funérailles — the funeral (formal)}

\label{lesson:FR-C314-funerailles}



\begin{quote}
Before the new one: say the French for to bury, then the French for the burial.
\end{quote}

\subsection*{You'll want to know: les funérailles}
\textbf{les funérailles} — \textquotedblleft{}the funeral (formal)\textquotedblright{}.

\begin{quote}
Les funérailles du chanteur ont eu lieu dans son village.
\end{quote}

\begin{quote}
\emph{The singer's funeral was held in his village.}
\end{quote}

Formal and plural only; the everyday word is \emph{l'enterrement}.

\begin{cousinweb}
\emph{funérailles} comes from Medieval Latin \emph{fūnerālia} (EtymDB), built on \emph{fūnus} \textquotedblleft{}a funeral\textquotedblright{} (Lewis \& Short).
\end{cousinweb}

\begin{grammarlens}[title={What you've built}]
One more word for telling, explaining and describing at B1.
\end{grammarlens}

\subsection*{Your turn}
\begin{itemize}
\raggedright
  \item \emph{Say it:} \emph{les funérailles}
  \item \emph{Say it:} \emph{les funérailles}, once more
  \item \emph{Say it:} a sentence of your own with \emph{les funérailles}
  \item \emph{Recall:} say the French for to bury, then the French for the burial, then say \emph{les funérailles} again
\end{itemize}

\begin{tcolorbox}[breakable,colback=teal!4,colframe=teal!35!black,title={Before you move on}]
What does les funérailles mean? (\textquotedblleft{}The funeral (formal)\textquotedblright{}.) Say it once more.
\end{tcolorbox}
\section[le cimetière]{le cimetière — the cemetery, the graveyard}

\label{lesson:FR-C314-cimetiere}



\begin{quote}
Before the new one: say the French for the burial, then the French for the funeral.
\end{quote}

\subsection*{You'll want to know: le cimetière}
\textbf{le cimetière} — \textquotedblleft{}the cemetery, the graveyard\textquotedblright{}.

\begin{quote}
Le cimetière est derrière l'église.
\end{quote}

\begin{quote}
\emph{The cemetery is behind the church.}
\end{quote}

Masculine: \emph{le cimetière}.

\begin{cousinweb}
\emph{cimetière} comes from Church Latin \emph{coemētērium}, a borrowing from Greek \emph{koimēt\'{\={e}}rion} \textquotedblleft{}a sleeping place\textquotedblright{} (Lewis \& Short).
\end{cousinweb}

\begin{grammarlens}[title={What you've built}]
One more word for telling, explaining and describing at B1.
\end{grammarlens}

\subsection*{Your turn}
\begin{itemize}
\raggedright
  \item \emph{Say it:} \emph{le cimetière}
  \item \emph{Say it:} \emph{le cimetière}, once more
  \item \emph{Say it:} a sentence of your own with \emph{le cimetière}
  \item \emph{Recall:} say the French for the burial, then the French for the funeral, then say \emph{le cimetière} again
\end{itemize}

\begin{tcolorbox}[breakable,colback=teal!4,colframe=teal!35!black,title={Before you move on}]
What does le cimetière mean? (\textquotedblleft{}The cemetery, the graveyard\textquotedblright{}.) Say it once more.
\end{tcolorbox}
